\chapter{Proveniência e procedimentos do \textit{linkage} de origem}\label{ap:linkage-origem}

O \textit{linkage} (vinculação de registros) de origem antecede a avaliação desta tese. A epidemiologista Daniela Medina Coeli Bastos, do IESC-UFRJ, executou sozinha o \textit{linkage} no OpenRecLink (versão original, em C), o Comparador de Registros (CdR) e a revisão clerical que originou PAR. O autor não participou dessas etapas e recebeu o extrato pós-CdR (\path{data/COMPARADORSEMIDENT.csv}) como membro do LinkDataPop. Houve revisora única, sem dupla classificação independente. Conforme a informação de proveniência confirmada pelo autor e registrada em \path{data/proveniencia_oliveira_meta.json}, o CdR em Python foi executado antes da revisão, com nota e subescores disponíveis à revisora. Não está documentado o que ela efetivamente consultou durante a classificação, inclusive eventual consulta ao SIM, ao Sinan ou ao endereço. Essa distinção não elimina o potencial de viés de incorporação.

O extrato provém do processo de \textit{linkage} descrito por \citeonline{Oliveira2026obitoTBpulmonar}, entre os óbitos de residentes no município do Rio de Janeiro com tuberculose como causa básica (SIM, 2019 a 2022) e as notificações de tuberculose (Sinan, 2013 a 2022), em quatro etapas: \textit{linkage} determinístico, 29 passos probabilísticos no OpenRecLink, um passo aberto e deduplicação (p.~33 a~39). Os vínculos encontrados na etapa determinística foram retirados antes do \textit{linkage} probabilístico, de modo que os potenciais pares do extrato correspondem ao que restou dessa etapa. O extrato reúne os potenciais pares de três dos 29 passos probabilísticos (6, 9 e 11), cujos parâmetros, pontos de corte (7,0240; 7,8444; 7,9019) e contagens de vínculos pelo corte e após a revisão coincidem com os descritos na monografia (Quadro~5, p.~36 a~37). Esses três passos contribuíram com 247 dos 2.549 pares do processo completo (p.~35); a proporção de 1:249 refere-se, portanto, à base exportada. A correspondência documental e as decisões de proveniência informadas pelo autor estão registradas em \texttt{data/proveniencia\_oliveira\_meta.json}. O órgão fornecedor das bases, a data da extração e o critério de escolha dos três passos exportados não constam da documentação disponível ao autor. A escolha dessas bases justifica-se pelas razões detalhadas na Seção \ref{sec:justificativas-especificas}, em particular a relevância da tuberculose como condição marcadora da qualidade do cuidado em saúde e a experiência acumulada do grupo de pesquisa no \textit{linkage} dessas fontes de dados \cite{Oliveira2012uso, Bartholomay2014improved, Oliveira2016accuracy}.

O procedimento de \textit{linkage} probabilístico no OpenRecLink seguiu o modelo de decisão de Fellegi e Sunter \citeyearpar{Fellegi1969theory}. Na etapa de padronização, os campos textuais de ambas as bases foram convertidos para caixa alta, com remoção de acentos diacríticos, preposições e artigos, e as datas de nascimento foram normalizadas para formato numérico contínuo (AAAAMMDD). A etapa de comparação no OpenRecLink utilizou três campos de identificação: nome próprio e nome da mãe, comparados por algoritmo aproximado (tipo APROX, com limiar de similaridade de 85\%); e data de nascimento, comparada caractere a caractere (tipo CARACTERE, com limiar de 65\%). A variável sexo não integrou a etapa de comparação, tendo sido utilizada exclusivamente como chave de blocagem em passos específicos. Os parâmetros de concordância e discordância atribuídos a cada campo estão detalhados na Tabela~\ref{tab:parametros-orl}. Para cada potencial par, o programa calculou um escore probabilístico correspondente à soma das razões de log-verossimilhança (\textit{log-likelihood ratios}) dos campos comparados, no qual campos concordantes contribuem positivamente e campos discordantes contribuem negativamente para o escore total \cite{Fellegi1969theory, Christen2012book}. Na base utilizada, os escores probabilísticos do OpenRecLink variaram de $-12{,}85$ (discordância em todos os campos) a $+15{,}58$ (concordância plena), distribuindo-se em 116 valores distintos. Para evitar ambiguidade, o presente trabalho distingue duas grandezas: o escore do \textit{linkage} probabilístico, calculado pelo OpenRecLink como soma das razões de log-verossimilhança, e a nota final do Comparador de Registros (CdR; Seção~\ref{sec:comparador-registros}), escore composto derivado dos 29 subescores de similaridade campo a campo. Na origem, a revisão manual incidiu sobre a região abaixo do corte do escore OpenRecLink \cite{Oliveira2026obitoTBpulmonar}; nos resultados desta tese, área cinza da nota designa o intervalo da nota final do CdR em $[5,9)$ (Seção~\ref{sec:comparador-registros}). As duas regiões não coincidem.

\begin{table}[!ht]
\centering
\caption{Parâmetros de comparação configurados no OpenRecLink para os campos de identificação.}
\label{tab:parametros-orl}
\begin{tabular}{
    p{\dimexpr0.32\linewidth-2\tabcolsep\relax}
    >{\centering\arraybackslash}p{\dimexpr0.16\linewidth-2\tabcolsep\relax}
    *{2}{>{\centering\arraybackslash}p{\dimexpr0.20\linewidth-2\tabcolsep\relax}}
    >{\centering\arraybackslash}p{\dimexpr0.12\linewidth-2\tabcolsep\relax}}
\hline
\textbf{Campo} & \textbf{Tipo} & \textbf{Concordância (\%)} & \textbf{Discordância (\%)} & \textbf{Limiar (\%)} \\
\hline
Nome (NOME) & Aproximado & 99 & 1 & 85 \\
Nome da mãe (NOMEMAE) & Aproximado & 75 & 2 & 85 \\
Data de nascimento (DTNASC) & Caractere & 95 & 5 & 65 \\
\hline
\end{tabular}
\end{table}

A geração de potenciais pares no OpenRecLink empregou a técnica de blocagem (\textit{blocking}), que particiona os registros em subgrupos definidos por chaves comuns e restringe as comparações aos pares que compartilham a mesma chave, reduzindo substancialmente o custo computacional \cite{Christen2012book}. No OpenRecLink, os parâmetros de blocagem foram definidos por combinações de variáveis fonéticas e demográficas (código \textit{Soundex} do primeiro e do último bloco do nome, sexo, ano de nascimento, primeiro e último bloco do nome da mãe, código do município de residência e bairro de residência (BAIRES)), organizadas em passos sequenciais. Adotou-se, ainda, a estratégia de exclusão progressiva dos registros do Sinan-TB já pareados em cada passo, de modo que, uma vez identificado um potencial par, o respectivo registro do Sinan não retornasse ao universo de busca dos passos subsequentes. Foram utilizados múltiplos passos de blocagem (\textit{blocking steps}), estratégia recomendada para maximizar a sensibilidade do processo ao permitir que pares não capturados por uma chave restritiva fossem recuperados por chaves alternativas \cite{Coeli2002blocking}. O procedimento incluiu 29 passos probabilísticos e um passo aberto, cuja chave de blocagem foi o \textit{Soundex} do primeiro fragmento do nome do paciente (PBLOCO), conforme a descrição metodológica e o Quadro~6 de \citeonline{Oliveira2026obitoTBpulmonar} (p.~34 e~39). Três dos passos probabilísticos integram o extrato, sem informação sobre o critério de escolha para exportação. Os três passos exportados, com suas respectivas chaves de blocagem, foram:

\begin{enumerate}
\item \textbf{Passo por nome e nome da mãe (passo~6):} código \textit{Soundex} do primeiro e do último bloco do nome (PBLOCO, UBLOCO) e do primeiro e do último bloco do nome da mãe (PBLOCOMAE, UBLOCOMAE), gerando 2.649 potenciais pares, dos quais 80 foram classificados como pares (76 pelo ponto de corte e 4 em revisão manual);
\item \textbf{Passo restritivo (passo~9):} código \textit{Soundex} do primeiro e do último bloco do nome (PBLOCO, UBLOCO) + ano de nascimento (ANONASC) + sexo (SEXO), gerando 1.449 potenciais pares, dos quais 127 foram classificados como pares (120 pelo ponto de corte e 7 em revisão manual);
\item \textbf{Passo ampliado (passo~11):} código \textit{Soundex} do primeiro e do último bloco do nome (PBLOCO, UBLOCO) + sexo (SEXO) + código do município de residência (CODMUNRES), gerando 57.598 potenciais pares, dos quais 40 foram classificados como pares (15 pelo ponto de corte e 25 em revisão manual).
\end{enumerate}

O Quadro~5 de \citeonline{Oliveira2026obitoTBpulmonar} (p.~36 a~38) registra cortes de 7,0240, 7,8444 e 7,9019 para os passos 6, 9 e 11. As contagens pelo corte são 76, 120 e 15, e, após revisão, 80, 127 e 40, respectivamente; as diferenças correspondem a 4, 7 e 25 vínculos PAR igual a 2 (\path{data/par_por_passo.csv}). A coincidência até a quarta casa decimal com o menor escore entre PAR igual a 1 no extrato decorre da definição do \textit{label}, não de validação independente. Não se documentaram a ordenação dos candidatos nem a regra de parada da revisão.

Na origem, o CdR foi usado na revisão do passo aberto, que recuperou 5 pares \cite{Oliveira2026obitoTBpulmonar} (p.~39); o cálculo da nota final sobre os potenciais pares desta tese foi outra execução, com a implementação em Python.

O texto da qualificação, preservado em \path{qualificacao-2026:text/latex/capitulo5.tex}, atribuía a Luís Guilherme Buteri Alves o pré-processamento, incluindo limpeza de nomes inconsistentes, como ``homem pardo'', ``morador de rua'' e ``recém-nascido'', e criação de DTNASC2 para normalização das datas de nascimento. Essa é informação da qualificação, não documentada em fonte primária de execução acessível ao autor; sua recuperação preserva o histórico do relato, sem comprovar os procedimentos. Essa atribuição ao pré-processamento não corresponde a uma segunda revisão clerical nem altera a autoria única do \textit{linkage}, do CdR e da revisão confirmada acima. As afirmações da qualificação sobre ordenação decrescente e consulta ao SIM, ao Sinan e ao endereço não são retomadas como fatos, pois não dispõem de documentação independente.
